\documentclass[11pt]{article}
\usepackage[margin=1in]{geometry}
\usepackage{amsmath,amssymb,amsthm}
\usepackage{enumitem}

\newcommand{\Form}{\mathrm{Form}}
\newcommand{\T}{\mathrm{T}}
\newcommand{\F}{\mathrm{F}}
\newcommand{\TF}{\{\T,\F\}}
\newcommand{\ext}[1]{\overline{#1}}
\newcommand{\Sfive}{\{\lnot,\land,\lor,\to,\leftrightarrow\}}
\newcommand{\vars}{\mathrm{vars}}
\newcommand{\bin}{\mathbin{\circ}}

\newcommand{\solpart}[1]{\smallskip\noindent\textbf{(#1)}\ }

\newcommand{\prob}[2]{\subsection*{Problem #1 --- #2}}

\begin{document}

\begin{center}
{\LARGE Solutions}\\[3pt]
{\large Problem Set: The Interpretation of Propositional Formulas}\\[2pt]
{\large Goldrei \S 2.3, with emphasis on Theorem 2.2}
\end{center}

\bigskip

\noindent\emph{How to read these.} Throughout, ``Theorem 2.2'' refers to the
existence-and-uniqueness of the extension $\ext{v}$; when a step uses only one
half I say which. ``Relevance'' means Problem 8(a). ``Unique readability'' is the
\S 2.2 fact that every formula of length $\ge 1$ has a unique principal connective
and hence a unique immediate decomposition. I write $\vars(\varphi)$ for the set
of propositional variables occurring in $\varphi$, and for a truth assignment $v$
I use the truth-table operations on values freely (e.g.\ $v(\lnot\theta)=\lnot
v(\theta)$). Many later problems cite earlier ones, as the ground rules allow.

\subsection*{Verdict summary (full arguments below)}

\begin{tabbing}
\hspace{3.3cm}\=\kill
Problem 1: \> (a)~T \quad(b)~T \quad(c)~F \quad(d)~T \quad(e)~T \quad(f)~F \quad(g)~T\\[2pt]
Problem 6(e): \> $((p\to q)\lor(q\to p))$ taut.; \ $(p\leftrightarrow\lnot p)$ contra.; \ $((p\lor q)\to p)$ neither; \ De Morgan taut.\\[2pt]
Problem 7: \> (a)~$2^{n-2}$ \quad(b)~$2^{n-3}$ \quad(c)~$2^{n-2}$ \quad(d)~$3$\\[2pt]
Problem 11: \> (a)~F \quad(b)~F \quad(c)~T \quad(d)~F \quad(e)~T \quad(f)~F \quad(g)~T \quad(h)~T\\[2pt]
Problem 12: \> (a)~T \quad(b)~T \quad(c)~F \quad(d)~T \quad(e)~T \quad(f)~F\\[2pt]
Problem 13: \> (a)~I \quad(b)~III \quad(c)~I \quad(d)~II \quad(e)~II \quad(f)~III \quad(g)~existence half\\
\end{tabbing}

\section*{Part I --- Reading Theorem 2.2}

\prob{1}{Reading the theorem precisely}

Assume $P\neq\varnothing$ (if $P=\varnothing$ the language is empty and every
statement degenerates); fix a variable $p_0\in P$ for use in refutations.

\solpart{a} \textbf{True.} This is exactly the existence half of Theorem 2.2:
$w=\ext{v}$ is a truth assignment with $\ext{v}(p)=v(p)$ for all $p\in P$.

\solpart{b} \textbf{True.} This is the uniqueness half of Theorem 2.2: at most one
truth assignment agrees with $v$ on $P$. (Together with (a): exactly one.)

\solpart{c} \textbf{False.} The condition constrains $w$ only on the variables in
$P$; on every compound formula $w$ is unconstrained. Take $w_1,w_2\colon
\Form(P,S)\to\TF$ both agreeing with $v$ on every variable, but with
$w_1(\lnot p_0)=\T$ and $w_2(\lnot p_0)=\F$ (define them arbitrarily, but
differently, here; e.g.\ set each equal to $\ext{v}$ except at $\lnot p_0$). Both
satisfy ``$w(p)=v(p)$ for all $p\in P$'' since $\lnot p_0$ is not a variable, and
$w_1\ne w_2$. So there is more than one such function; ``exactly one'' fails.

\solpart{d} \textbf{True.} \emph{Existence of $v$:} put $v=w{\restriction}P$
(the restriction of $w$ to the variables). Then $w$ is a truth assignment agreeing
with $v$ on $P$, so by the uniqueness half of Theorem 2.2, $w=\ext{v}$.
\emph{Uniqueness of $v$:} if $\ext{v}=w=\ext{v'}$ then for each $p\in P$,
$v(p)=\ext{v}(p)=w(p)=\ext{v'}(p)=v'(p)$, so $v=v'$.

\solpart{e} \textbf{True.} Both $u$ and $w$ are truth assignments extending the
common function $v:=u{\restriction}P=w{\restriction}P$. By the uniqueness half of
Theorem 2.2 there is only one such extension, so $u=w$; in particular
$u(\varphi)=w(\varphi)$ for every $\varphi$. (This is the book's own derivation of
``a truth assignment is determined by its values on the variables''; it is also
Problem 8(a) taken over the full variable set.)

\solpart{f} \textbf{False.} Agreeing with a truth assignment on the variables does
not force $w$ to respect the truth tables at compound formulas. Let $u=\ext{v}$
where $v(p_0)=\T$, so $u(\lnot p_0)=\F$. Define $w$ to equal $u$ everywhere except
$w(\lnot p_0)=\T$. Then $w$ is a genuine function $\Form(P,S)\to\TF$ agreeing with
$u$ on every variable (they agree at $p_0$ and at all other variables). But $w$
violates the $\lnot$-clause at $\lnot p_0$: respecting $\lnot$ would require
$w(\lnot p_0)=\lnot w(p_0)=\lnot\T=\F$, whereas $w(\lnot p_0)=\T$. So $w$ is not a
truth assignment.

\solpart{g} \textbf{True.} For each variable $p$, the formula $\lnot p$ has length
$1$, so by hypothesis $u(\lnot p)=w(\lnot p)$. Since $u,w$ respect $\lnot$, this
says $\lnot u(p)=\lnot w(p)$, hence $u(p)=w(p)$. Thus $u$ and $w$ agree on every
variable, so by (e) (equivalently the uniqueness half of Theorem 2.2) $u=w$.

\prob{2}{Extensions that are not truth assignments}

Here $P=\{p\}$, $v(p)=\T$.

\solpart{a} Let $f$ be the constant function $\T$ (so $f(\chi)=\T$ for every
$\chi$). Then $f(p)=\T$, but $f$ violates the $\lnot$-clause at $\lnot p$:
$f(\lnot p)=\T$ while respecting $\lnot$ would demand $\lnot f(p)=\F$.
Let $g$ agree with the genuine assignment $\ext{v}$ everywhere except
$g\big((p\land p)\big)=\F$. Then $g(p)=\ext{v}(p)=\T$, but $g$ violates the
$\land$-clause at $(p\land p)$: since $g(p)=\T$, respecting $\land$ would require
$g\big((p\land p)\big)=\T\land\T=\T$, whereas it is $\F$. Finally $f\neq g$: they
disagree at $\lnot p$ (indeed $g(\lnot p)=\ext{v}(\lnot p)=\F\neq\T=f(\lnot p)$).

\solpart{b} $\Form(\{p\},S)$ is infinite: $p,\lnot p,\lnot\lnot p,\lnot\lnot\lnot
p,\dots$ are pairwise distinct formulas. For each $k\ge 1$ define
$f_k\colon\Form(\{p\},S)\to\TF$ by $f_k(p)=\T$, $f_k(\lnot^{\,k}p)=\T$, and
$f_k(\chi)=\F$ for all other $\chi$. Each $f_k$ extends $v$ (value $\T$ at $p$),
and $f_j\ne f_k$ for $j\ne k$ (they differ at $\lnot^{\,j}p$). So there are
infinitely many functions extending $v$. (Almost none are truth assignments; the
problem asks only for functions.)

\solpart{c} \textbf{Exactly one}, by the uniqueness half of Theorem 2.2. The
freedom exhibited in (b) is freedom among arbitrary functions; requiring the
extension to respect the truth tables collapses it to a single assignment.

\solpart{d} Let $\mathcal{A}$ be the set of truth assignments on $\Form(P,S)$ and
$\mathcal{V}$ the set of functions $P\to\TF$. Define
$\rho\colon\mathcal{A}\to\mathcal{V}$ by $\rho(w)=w{\restriction}P$ and
$E\colon\mathcal{V}\to\mathcal{A}$ by $E(v)=\ext{v}$ (well-defined into
$\mathcal{A}$ by the existence half of Theorem 2.2).
\emph{$\rho\circ E=\mathrm{id}$:} $\rho(E(v))=\ext{v}{\restriction}P=v$, since
$\ext{v}(p)=v(p)$ for all $p\in P$.
\emph{$E\circ\rho=\mathrm{id}$:} $E(\rho(w))=\overline{w{\restriction}P}$; but $w$
is a truth assignment extending $w{\restriction}P$, so by the uniqueness half of
Theorem 2.2 the unique such extension is $w$ itself, i.e.\
$\overline{w{\restriction}P}=w$.
Being mutually inverse, $\rho$ and $E$ are bijections. Hence if $|P|=n$ then
$|\mathcal{A}|=|\mathcal{V}|=|\TF|^{\,|P|}=2^n$ (cf.\ Exercise 2.19).

\prob{3}{What unique readability is for}

$\sigma$ is the bracket-free string $p\land q\lor r$.

\solpart{a} Two construction trees:
\[
T_1:\ (p\land q)\lor r
\qquad\text{(root }\lor\text{, children }p\land q\text{ and }r),
\]
\[
T_2:\ p\land(q\lor r)
\qquad\text{(root }\land\text{, children }p\text{ and }q\lor r).
\]
As \emph{strings} both spell $p\land q\lor r$: concatenating $p\land q$, $\lor$,
$r$ gives $\sigma$, and so does concatenating $p$, $\land$, $q\lor r$. Same string,
two different trees.

\solpart{b} With $v(p)=\F$, $v(q)=v(r)=\T$:
\[
T_1:\quad p\land q\mapsto \F\land\T=\F,\qquad \F\lor r\mapsto \F\lor\T=\T;
\]
\[
T_2:\quad q\lor r\mapsto \T\lor\T=\T,\qquad p\land\T\mapsto \F\land\T=\F.
\]
So the same string receives value $\T$ along $T_1$ and $\F$ along $T_2$. The
recursive clauses therefore do not determine a single value for $\sigma$: there is
no well-defined function from pseudo-formulas to $\TF$ produced by these clauses.

\solpart{c} The failing sentence is the one asserting the principal connective in
the recursive step: Goldrei writes that a formula of length $n+1$ ``has a principal
connective \emph{(which is unique)}, so is one of the forms $\lnot\theta$,
$(\theta\land\psi),\dots$ where $\theta$ and $\psi$ are of length $\le n$,'' and
then defines the value ``using the appropriate row of the truth table for its
principal connective.'' For pseudo-formulas this uniqueness fails: $\sigma$ has
\emph{two} principal connectives ($\lor$ via $T_1$, $\land$ via $T_2$) and two
decompositions into $(\theta,\psi)$. So ``the appropriate row for its principal
connective'' is not well specified --- there are two, giving two candidate values.
In genuine $\Form(P,S)$, unique readability guarantees exactly one principal
connective and one decomposition, so the recursive rule has a single output at
each step.

\solpart{d} What is threatened is the \textbf{well-definedness} of $\ext{v}$ as a
function. Producing \emph{a} value is not the issue (each tree yields one); and
uniqueness in the sense of Theorem 2.2 --- at most one \emph{truth assignment} ---
is a separate, downstream question. Part (b) shows the \emph{rule}
``$\ext{v}(\sigma):=$ the value from the decomposition of $\sigma$'' fails to
name one value, because $\sigma$ has more than one decomposition with differing
values. The three notions line up as a dependency chain: unique readability makes
the recursive clauses well-defined (each formula has a unique immediate
decomposition, so the rule has a unique output); well-definedness is what lets
$\ext{v}$ be a function at all; and only once $\ext{v}$ is a function can one ask
whether it exists as a truth assignment and whether it is the unique one. Removing
unique readability breaks the very first link, which is what (b) exhibits.

\prob{4}{Adding a connective}

\solpart{a} $v$ \emph{respects the truth table of $\oplus$} means: for all
$\theta,\psi\in\Form(P,S_\oplus)$,
\[
v\big((\theta\oplus\psi)\big)=
\begin{cases}
\T, & \text{if exactly one of }v(\theta),v(\psi)\text{ is }\T,\\
\F, & \text{if }v(\theta)=v(\psi).
\end{cases}
\]

\solpart{b} \emph{Analogue of Theorem 2.2.} Let $P$ be a set of propositional
variables, let $S_\oplus=\{\lnot,\land,\lor,\to,\leftrightarrow,\oplus\}$, and let
$v\colon P\to\TF$ be a function. Then there is a unique truth assignment
$\ext{v}\colon\Form(P,S_\oplus)\to\TF$ with $\ext{v}(p)=v(p)$ for all $p\in P$
(where ``truth assignment'' now means a function respecting the truth tables of
all six connectives in $S_\oplus$).

\solpart{c} \emph{Existence (recursion on length).} One added case in the
recursive step: if the principal connective is $\oplus$, so $\varphi=(\theta\oplus
\psi)$ with $\theta,\psi$ of length $\le n$, define $\ext{v}(\varphi)$ by the
$\oplus$-row determined by $\ext{v}(\theta),\ext{v}(\psi)$. Everything else is
unchanged. \emph{Uniqueness (induction on length).} One added case: if
$u,w$ are truth assignments (now respecting all six tables) agreeing on the
variables and $\varphi=(\theta\oplus\psi)$, then the induction hypothesis gives
$u(\theta)=w(\theta)$ and $u(\psi)=w(\psi)$, and since both respect the
$\oplus$-table, $u(\varphi)=w(\varphi)$. Otherwise unchanged.

\solpart{d} The facts needed about $\Form(P,S_\oplus)$, of the kind proved for
$\Form(P,S)$ in \S 2.2, are: (i) every formula of length $\ge 1$ has a
\emph{unique} principal connective, and correspondingly a unique immediate
decomposition --- either $\lnot\theta$ for a unique $\theta$, or $(\theta\bin\psi)$
for a unique connective $\bin\in\{\land,\lor,\to,\leftrightarrow,\oplus\}$ and
unique $\theta,\psi$ (unique readability for the enlarged language); and (ii) those
immediate subformulas have strictly smaller length. If uniqueness of the principal
connective failed, then in the existence proof the step ``define $\ext{v}(\varphi)$
by the table for its principal connective'' would be ambiguous exactly as in
Problem 3 --- a formula might decompose both as $(\theta\oplus\psi)$ and as
$(\theta'\lor\psi')$, giving two candidate values, so $\ext{v}$ would not be
well-defined; and in the uniqueness proof the case split on the principal
connective would be neither exhaustive nor exclusive, so the inductive step would
fail to pin down $u(\varphi)$.


\section*{Part II --- Computing with truth assignments}

\prob{5}{Bottom-up evaluation}

Here $w(p)=\F$, $w(q)=\T$, $w(r)=\T$.

\solpart{a} $w\big(((\lnot p\lor q)\to(r\land\lnot q))\big)$:
\[
\lnot p=\T,\quad (\lnot p\lor q)=\T\lor\T=\T,\quad \lnot q=\F,\quad
(r\land\lnot q)=\T\land\F=\F,
\]
\[
\big((\lnot p\lor q)\to(r\land\lnot q)\big)=\T\to\F=\boxed{\F}.
\]

\solpart{b} $w\big(((p\to(q\to r))\leftrightarrow\lnot(p\lor\lnot r))\big)$:
\[
(q\to r)=\T\to\T=\T,\quad (p\to(q\to r))=\F\to\T=\T,
\]
\[
\lnot r=\F,\quad (p\lor\lnot r)=\F\lor\F=\F,\quad \lnot(p\lor\lnot r)=\T,
\]
\[
\big((p\to(q\to r))\leftrightarrow\lnot(p\lor\lnot r)\big)=\T\leftrightarrow\T
=\boxed{\T}.
\]

\solpart{c} \textbf{Determined; value $\T$.} We have
$((q\land s)\lor r)$ with $w(r)=\T$, and $x\lor\T=\T$ for both values of $x$.
So whatever $w(s)$ is, $(q\land s)\lor r=\T$. Every $w$ consistent with the data
gives $\T$.

\solpart{d} \textbf{Not determined.} Here $(q\land s)\land r$; with $w(q)=w(r)=\T$
this equals $(\T\land w(s))\land\T=w(s)$, which is unconstrained. Two truth
assignments consistent with the data:
\[
w_1:\ (p,q,r,s)\mapsto(\F,\T,\T,\T)\ \Rightarrow\ (q\land s)\land r=(\T\land\T)\land\T=\T,
\]
\[
w_2:\ (p,q,r,s)\mapsto(\F,\T,\T,\F)\ \Rightarrow\ (q\land s)\land r=(\T\land\F)\land\T=\F.
\]
Both exist as truth assignments by the existence half of Theorem 2.2 (they are
specified by consistent values on the variables), and they give different values.

\solpart{e} Only $w(p)=\T$ is known.
\begin{itemize}[nosep]
\item $(p\lor q)$: $\T\lor w(q)=\T$ always. \textbf{Determined}, value $\T$.
\item $(p\land q)$: $\T\land w(q)=w(q)$, unconstrained. \textbf{Not determined}
($w(q)=\T$ gives $\T$; $w(q)=\F$ gives $\F$).
\item $(q\to p)$: $w(q)\to\T=\T$ always. \textbf{Determined}, value $\T$.
\item $(q\lor\lnot q)$: $w(q)\lor\lnot w(q)=\T$ always. \textbf{Determined}, value
$\T$ (indeed independent of $p$ as well).
\item $(q\land\lnot q)$: $w(q)\land\lnot w(q)=\F$ always. \textbf{Determined},
value $\F$.
\end{itemize}

\prob{6}{Truth tables, concrete and schematic}

\solpart{a} Table of $((p\lor\lnot q)\to(q\land r))$:
\[
\begin{array}{ccc|c|c|c}
p&q&r& p\lor\lnot q & q\land r & ((p\lor\lnot q)\to(q\land r))\\\hline
\T&\T&\T& \T&\T&\T\\
\T&\T&\F& \T&\F&\F\\
\T&\F&\T& \T&\F&\F\\
\T&\F&\F& \T&\F&\F\\
\F&\T&\T& \F&\T&\T\\
\F&\T&\F& \F&\F&\T\\
\F&\F&\T& \T&\F&\F\\
\F&\F&\F& \T&\F&\F
\end{array}
\]
(True in exactly the three rows $\T\T\T$, $\F\T\T$, $\F\T\F$; note the two
$\T$'s in rows $5,6$ come from a \emph{false} antecedent.)

\solpart{b} Table of $((\varphi\to\psi)\to\varphi)$ in the values of $\varphi,\psi$:
\[
\begin{array}{cc|c|c}
\varphi&\psi& (\varphi\to\psi) & ((\varphi\to\psi)\to\varphi)\\\hline
\T&\T& \T&\T\\
\T&\F& \F&\T\\
\F&\T& \T&\F\\
\F&\F& \T&\F
\end{array}
\]
The last column equals the $\varphi$-column in every row: $((\varphi\to\psi)\to
\varphi)$ takes the same value as $\varphi$ under every truth assignment.

\solpart{c} Realizing $(\varphi,\psi)=(\F,\T)$: take $\varphi_0=p$, $\psi_0=q$, and
$w$ with $w(p)=\F$, $w(q)=\T$. Then $w(\varphi_0)=\F$, $w(\psi_0)=\T$, and
\[
w\big(((p\to q)\to p)\big):\quad (p\to q)=\F\to\T=\T,\quad (\T\to p)=\T\to\F=\F.
\]
So the value is $\F$, matching the row $(\F,\T)$ of (b) (whose entry equals
$\varphi=\F$).

\solpart{d} Put $\psi:=\lnot\varphi$. Under any $w$, $w(\psi)=\lnot w(\varphi)$, so
the pair $(w(\varphi),w(\psi))$ is always $(\T,\F)$ or $(\F,\T)$: the rows $(\T,\T)$
and $(\F,\F)$ become \textbf{unrealizable}, since they would require $w(\varphi)=
w(\psi)=w(\varphi)=\lnot w(\varphi)$, impossible. On the surviving rows, and indeed
on all of them, (b) gives $((\varphi\to\psi)\to\varphi)=\varphi$; hence
$((\varphi\to\lnot\varphi)\to\varphi)$ takes the same value as $\varphi$ under every
$w$. Therefore it is a tautology iff $\varphi$ is true under every $w$, i.e.\ iff
$\varphi$ is a tautology; and a contradiction iff $\varphi$ is false under every
$w$, i.e.\ iff $\varphi$ is a contradiction.

\solpart{e} \emph{Classifications.}
\begin{itemize}[nosep]
\item $((p\to q)\lor(q\to p))$ --- \textbf{tautology.} If both disjuncts were
false under some $w$, then $(p\to q)=\F$ gives $w(p)=\T,w(q)=\F$, while $(q\to p)=
\F$ gives $w(q)=\T,w(p)=\F$; incompatible. So at least one disjunct is $\T$ under
every $w$.
\item $(p\leftrightarrow\lnot p)$ --- \textbf{contradiction.} $w(p)\leftrightarrow
\lnot w(p)$ is $\T\leftrightarrow\F=\F$ or $\F\leftrightarrow\T=\F$.
\item $((p\lor q)\to p)$ --- \textbf{neither.} $w(p)=\T$ gives $\T\to\T=\T$;
$w(p)=\F,w(q)=\T$ gives $\T\to\F=\F$.
\item $(\lnot(p\land q)\leftrightarrow(\lnot p\lor\lnot q))$ --- \textbf{tautology}
(De Morgan). Checking the four rows: $(\T,\T)\!:\F\!\leftrightarrow\!\F=\T$;
$(\T,\F)\!:\T\!\leftrightarrow\!\T=\T$; $(\F,\T)\!:\T\!\leftrightarrow\!\T=\T$;
$(\F,\F)\!:\T\!\leftrightarrow\!\T=\T$.
\end{itemize}

\prob{7}{Counting assignments}

Values on the ``free'' variables (those not named in the formula) may be anything,
contributing a factor $2^{(\text{number free})}$; the named variables are pinned.

\solpart{a} $(p_1\land\lnot p_2)$ is true iff $p_1=\T$ and $p_2=\F$; $p_3,\dots,
p_n$ free. Count $=2^{n-2}$.

\solpart{b} $\lnot((p_1\lor p_2)\lor p_3)$ is true iff $(p_1\lor p_2\lor p_3)=\F$
iff $p_1=p_2=p_3=\F$; the rest free. Count $=2^{n-3}$.

\solpart{c} $\lnot(p_1\to p_2)$ is true iff $(p_1\to p_2)=\F$ iff $p_1=\T,p_2=\F$;
the rest free. Count $=2^{n-2}$.

\solpart{d} $((p_1\lor p_2)\land p_3)$ with $n=3$: need $p_3=\T$ and $(p_1\lor p_2)
=\T$. Fixing $p_3=\T$, the condition $p_1\lor p_2=\T$ excludes only $p_1=p_2=\F$,
leaving $3$ of the $4$ patterns of $(p_1,p_2)$. Count $=3$.

\section*{Part III --- Inductions on the length of a formula}

\prob{8}{The relevance lemma}

\solpart{a} \textbf{Claim.} If $u,w$ are truth assignments and $u(q)=w(q)$ for
every $q\in\vars(\varphi)$, then $u(\varphi)=w(\varphi)$.

\emph{Proof by induction on the length of $\varphi$.}
\emph{Induction hypothesis:} for every formula $\theta$ of length $\le n$, if $u,w$
agree on every variable in $\vars(\theta)$, then $u(\theta)=w(\theta)$.

\emph{Base ($\varphi$ of length $0$).} Then $\varphi$ is a variable $p$, and
$\vars(\varphi)=\{p\}$; the hypothesis gives $u(p)=w(p)$, i.e.\ $u(\varphi)=
w(\varphi)$.

\emph{Inductive step ($\varphi$ of length $n+1$).} By unique readability $\varphi$
has one principal connective.
\begin{itemize}[nosep]
\item $\varphi=\lnot\theta$. Then $\vars(\varphi)=\vars(\theta)$, so $u,w$ agree on
$\vars(\theta)$; as $\theta$ has length $n$, the IH gives $u(\theta)=w(\theta)$.
Since $u,w$ \emph{respect $\lnot$} (used here),
$u(\lnot\theta)=\lnot u(\theta)=\lnot w(\theta)=w(\lnot\theta)$.
\item $\varphi=(\theta\land\psi)$. Then $\vars(\varphi)=\vars(\theta)\cup
\vars(\psi)$, so $u,w$ agree on each of $\vars(\theta),\vars(\psi)$; both have
length $\le n$, so the IH gives $u(\theta)=w(\theta)$ and $u(\psi)=w(\psi)$. Since
$u,w$ \emph{respect $\land$} (used here), $u((\theta\land\psi))=u(\theta)\land
u(\psi)=w(\theta)\land w(\psi)=w((\theta\land\psi))$.
\item $\varphi=(\theta\to\psi)$. As above $u,w$ agree on $\vars(\theta),
\vars(\psi)$; the IH gives $u(\theta)=w(\theta)$, $u(\psi)=w(\psi)$; since $u,w$
\emph{respect $\to$} (used here), the value $u((\theta\to\psi))$, read off the
$\to$-table from $(u(\theta),u(\psi))=(w(\theta),w(\psi))$, equals
$w((\theta\to\psi))$.
\item The cases $\lor,\leftrightarrow$ are identical, using the respective tables.
\end{itemize}
This closes the induction. \hfill$\square$

\solpart{b} \emph{Exercise 2.18.} Suppose $v,v'$ are truth assignments agreeing at
every variable except possibly $p$ (that is, $v(q)=v'(q)$ for all $q\neq p$), and
let $\varphi$ be a formula (over $\{\lnot,\land,\lor\}$, as in the book) in which
$p$ does not occur. Every $q\in\vars(\varphi)$ satisfies $q\neq p$, so $v,v'$ agree
on all of $\vars(\varphi)$; by (a), $v(\varphi)=v'(\varphi)$. The analogue for the
full set $S$ is immediate, since (a) was proved for $S$. (Indeed (a) is stronger:
the two assignments may differ at \emph{any} set of variables disjoint from
$\vars(\varphi)$, not just at a single $p$.)

\solpart{c} \emph{Soundness of truth tables.} Write $V=\vars(\varphi)\subseteq
\{p_1,\dots,p_n\}$.

\emph{The bridge.} Part (a) as stated compares two assignments on a common domain.
For soundness we need the mild variant: \emph{if $u$ is a truth assignment on
$\Form(P,S)$ and $u'$ a truth assignment on $\Form(P',S)$ with $\vars(\varphi)
\subseteq P\cap P'$, and $u(q)=u'(q)$ for every $q\in\vars(\varphi)$, then
$u(\varphi)=u'(\varphi)$.} The proof is \emph{verbatim} the induction in (a): at
each node the current subformula uses only variables in $\vars(\varphi)\subseteq
P\cap P'$, hence lies in both $\Form(P,S)$ and $\Form(P',S)$ and the IH applies;
the argument never uses that the two domains coincide, only that every subformula
of $\varphi$ lives in both. This is the point the problem flags: ``agree on the
variables of $\varphi$'' is the right hypothesis even across different domains.

\emph{Soundness.} ($\Rightarrow$) If $\varphi$ is a tautology --- true under every
truth assignment defined on any variable set containing $\vars(\varphi)$ --- then
in particular it is true under each of the $2^n$ assignments on
$\Form(\{p_1,\dots,p_n\},S)$.

($\Leftarrow$) Suppose all $2^n$ assignments on $\Form(\{p_1,\dots,p_n\},S)$ make
$\varphi$ true. Let $w$ be an arbitrary truth assignment on some $\Form(P,S)$ with
$V\subseteq P$. Define $g\colon\{p_1,\dots,p_n\}\to\TF$ by $g(p_i)=w(p_i)$ when
$p_i\in P$ and $g(p_i)=\T$ otherwise (the arbitrary choice affects only $p_i\notin
V$, since $p_i\in V\Rightarrow p_i\in P$). By the existence half of Theorem 2.2,
$\ext{g}$ is one of the $2^n$ assignments on $\Form(\{p_1,\dots,p_n\},S)$, so
$\ext{g}(\varphi)=\T$ by hypothesis. Now $\ext{g}$ and $w$ agree on every
$q\in V=\vars(\varphi)$ (for such $q=p_i\in P$ we have $\ext{g}(p_i)=g(p_i)=
w(p_i)$), so the cross-domain bridge gives $w(\varphi)=\ext{g}(\varphi)=\T$. As $w$
was arbitrary, $\varphi$ is a tautology.

(If some $p_i\notin V$, the $2^n$ assignments simply test each genuinely relevant
pattern $2^{\,n-|V|}$ times; the criterion is unaffected.)

\prob{9}{Negating the variables: Exercise 2.28 extended}

Recall $v^*$ is the assignment with $v^*(p)=\lnot v(p)$ on variables (a truth
assignment by Theorem 2.2), and $\varphi^*$ is defined by the stated recursion
(each variable negated, connectives unchanged).

\solpart{a} \textbf{Claim.} $v(\varphi)=v^*(\varphi^*)$ for every truth assignment
$v$ and formula $\varphi$. Fix $v$ (hence $v^*$) and induct on the length of
$\varphi$; IH: this holds for all formulas of length $\le n$.

\emph{Base ($\varphi=p$).} $\varphi^*=\lnot p$, so
\[
v^*(\varphi^*)=v^*(\lnot p)=\lnot v^*(p)=\lnot(\lnot v(p))=v(p)=v(\varphi),
\]
using that $v^*$ respects $\lnot$ and the definition $v^*(p)=\lnot v(p)$.

\emph{Step.}
\begin{itemize}[nosep]
\item $\varphi=\lnot\theta$: $\varphi^*=\lnot\theta^*$, so
$v^*(\varphi^*)=\lnot v^*(\theta^*)=\lnot v(\theta)=v(\lnot\theta)=v(\varphi)$,
by the $\lnot$-clause (for $v^*$ and for $v$) and the IH on $\theta$.
\item $\varphi=(\theta\to\psi)$: $\varphi^*=(\theta^*\to\psi^*)$. Both
$v^*(\varphi^*)$ and $v(\varphi)$ are read off the \emph{same} $\to$-table, the
former from $(v^*(\theta^*),v^*(\psi^*))$ and the latter from $(v(\theta),v(\psi))$;
by the IH these input pairs coincide, so the outputs coincide.
\item $\varphi=(\theta\leftrightarrow\psi)$: identical, using the
$\leftrightarrow$-table.
\item $\land,\lor$ similar.
\end{itemize}
(The key structural point: $*$ negates variables but leaves each binary connective
\emph{unchanged}, so at each step the two sides use the identical truth table.)
\hfill$\square$

\solpart{b} $(v^*)^*$ is the assignment with $(v^*)^*(p)=\lnot v^*(p)=\lnot(\lnot
v(p))=v(p)$ for every variable $p$. So $(v^*)^*$ and $v$ are truth assignments
agreeing on all variables; by the uniqueness half of Theorem 2.2 they are equal as
functions on $\Form(P,S)$, i.e.\ $(v^*)^*=v$. Hence $v\mapsto v^*$ is its own
inverse, so a bijection of the set of truth assignments onto itself.

\solpart{c} \textbf{$\varphi$ tautology $\iff\varphi^*$ tautology.}
($\Rightarrow$) Let $w$ be an arbitrary truth assignment; we show $w(\varphi^*)=\T$.
Put $v:=w^*$ (a truth assignment). By (a), $v(\varphi)=v^*(\varphi^*)$, and by (b)
$v^*=(w^*)^*=w$; so $w(\varphi^*)=v(\varphi)$. Since $\varphi$ is a tautology,
$v(\varphi)=\T$, hence $w(\varphi^*)=\T$. As $w$ was arbitrary, $\varphi^*$ is a
tautology.
($\Leftarrow$) Let $w$ be arbitrary; by (a) with $v:=w$, $w(\varphi)=w^*(\varphi^*)$;
since $\varphi^*$ is a tautology, $w^*(\varphi^*)=\T$, so $w(\varphi)=\T$.

The quantifier point the problem asks for: in ($\Rightarrow$) we need
$w(\varphi^*)$ for \emph{arbitrary} $w$, and (a) only directly gives values of the
form $v^*(\varphi^*)$. This suffices precisely because $*$ is \emph{onto} (indeed a
bijection, by (b)): every assignment $w$ equals $v^*$ for some $v$ --- namely
$v=w^*$ --- so the universally quantified identity in (a) reaches every $w$.

\solpart{d} \textbf{$\varphi$ contradiction $\iff\varphi^*$ contradiction}, no new
induction. By Exercise 2.26, $\varphi$ is a contradiction iff $\lnot\varphi$ is a
tautology. Also $(\lnot\varphi)^*=\lnot(\varphi^*)$ by the definition of $*$.
Therefore
\[
\begin{aligned}
\varphi\text{ contradiction}&\iff\lnot\varphi\text{ tautology}
\overset{(c)}{\iff}(\lnot\varphi)^*\text{ tautology}\\
&\iff\lnot(\varphi^*)\text{ tautology}\iff\varphi^*\text{ contradiction},
\end{aligned}
\]
the middle step applying (c) to $\lnot\varphi$ and the last again by Exercise 2.26.

\solpart{e} \textbf{True:} $v(\varphi^*)=v^*(\varphi)$ for all $v,\varphi$, with no
new induction. Apply (a) to the assignment $u:=v^*$: $u(\varphi)=u^*(\varphi^*)$,
i.e.\ $v^*(\varphi)=(v^*)^*(\varphi^*)=v(\varphi^*)$, using $(v^*)^*=v$ from (b).
Reading the identity backwards, $v(\varphi^*)=v^*(\varphi)$.

\prob{10}{$\star$ The substitution lemma}

\solpart{a} By unique readability every formula falls under exactly one clause of
the recursion: it is the variable $p$, or a variable $q\neq p$, or $\lnot\theta$
for a unique $\theta$, or $(\theta\bin\psi)$ for a unique connective and unique
$\theta,\psi$. Since the clauses never compete, $\varphi[\chi/p]$ is specified
unambiguously --- the same definition-by-recursion justified by the existence half
of Theorem 2.2.

\solpart{b} \textbf{Claim.} $v(\varphi[\chi/p])=u(\varphi)$, where $u$ is the truth
assignment with $u(p)=v(\chi)$ and $u(q)=v(q)$ for $q\neq p$. (Such $u$ exists and
is unique by Theorem 2.2, applied to the function on variables sending $p\mapsto
v(\chi)$ and $q\mapsto v(q)$.) Fix $v$ (hence $u$) and induct on the length of
$\varphi$; IH: the identity holds for formulas of length $\le n$.

\emph{Base, two subcases.}
\begin{itemize}[nosep]
\item $\varphi=p$: $\varphi[\chi/p]=\chi$, so $v(\varphi[\chi/p])=v(\chi)=u(p)=
u(\varphi)$.
\item $\varphi=q$, $q\neq p$: $\varphi[\chi/p]=q$, so $v(\varphi[\chi/p])=v(q)=
u(q)=u(\varphi)$.
\end{itemize}

\emph{Step.}
\begin{itemize}[nosep]
\item $\varphi=\lnot\theta$: $\varphi[\chi/p]=\lnot(\theta[\chi/p])$, so
$v(\varphi[\chi/p])=\lnot v(\theta[\chi/p])=\lnot u(\theta)=u(\lnot\theta)=
u(\varphi)$, by the $\lnot$-clause and the IH.
\item $\varphi=(\theta\to\psi)$: then $\varphi[\chi/p]=(\theta[\chi/p]\to
\psi[\chi/p])$, and by the IH applied to $\theta$ and $\psi$,
\[
\big(v(\theta[\chi/p]),\,v(\psi[\chi/p])\big)=\big(u(\theta),\,u(\psi)\big).
\]
Reading both $v(\varphi[\chi/p])$ and $u((\theta\to\psi))$ off the $\to$-table from
this common pair gives $v(\varphi[\chi/p])=u((\theta\to\psi))=u(\varphi)$.
\item $\land,\lor,\leftrightarrow$ similar.
\end{itemize}
\hfill$\square$

\solpart{c} If $\varphi$ is a tautology then, for every truth assignment $v$,
$v(\varphi[\chi/p])=u(\varphi)=\T$ by (b) (with $u$ the assignment it names, itself
a genuine truth assignment). As $v$ was arbitrary, $\varphi[\chi/p]$ is a
tautology.

\solpart{d} \textbf{The converse fails.} Take $\varphi=(p\lor q)$, substitute for
$p$ the formula $\chi=\lnot q$. Then
\[
\varphi[\chi/p]=(\lnot q\lor q),
\]
which is a tautology: $w(\lnot q\lor q)=\lnot w(q)\lor w(q)=\T$ for every $w$. But
$\varphi=(p\lor q)$ is \emph{not} a tautology: the assignment $w$ with $w(p)=\F$,
$w(q)=\F$ gives $w((p\lor q))=\F\lor\F=\F$. So $\varphi[\chi/p]$ is a tautology
while $\varphi$ is not.

\solpart{e} By (b), the value of a formula built from $\varphi,\psi$ under $w$
equals the value of the corresponding two-variable formula under the assignment
sending $p\mapsto w(\varphi)$, $q\mapsto w(\psi)$; when $\varphi,\psi$ are distinct
variables, Theorem 2.2 lets $(w(\varphi),w(\psi))$ be set to any of the four pairs
independently, so every row of the schematic table is realized. When $\varphi,\psi$
are correlated --- e.g.\ $\psi=\lnot\varphi$ as in Problem 6(d) --- the pair is
constrained (there $w(\psi)=\lnot w(\varphi)$), so rows violating the constraint
have no realizing assignment and drop out. In short, distinct metavariables need
not be independently settable; only genuinely independent variables reach all rows.

\section*{Part IV --- Tautologies, contradictions, counterexamples}

\prob{11}{Calibration: true or false?}

\solpart{a} \textbf{False.} Agreement at $\varphi$ does not force agreement at its
variables. Take $\varphi=(p\lor q)$; let $u(p)=\T,u(q)=\T$ and $w(p)=\T,w(q)=\F$.
Then $u(\varphi)=\T$ and $w(\varphi)=\T\lor\F=\T$, so $u(\varphi)=w(\varphi)$, yet
$q\in\vars(\varphi)$ and $u(q)=\T\neq\F=w(q)$. (This is the \emph{converse} of
relevance, and it is false.)

\solpart{b} \textbf{False.} Any truth assignment respects $\land$:
$v((p\land q))=v(p)\land v(q)$. If $v(q)=\F$ then $v(p)\land\F=\F\neq\T$. So no
truth assignment does this.

\solpart{c} \textbf{True.} A mere function need not respect $\land$. Define $f$ with
$f((p\land q))=\T$, $f(q)=\F$, and $f$ arbitrary (say $\T$) elsewhere. This is a
well-defined function with the two required values. (Contrast (b): the $\land$-clause
is exactly what a function may flout but an assignment may not.)

\solpart{d} \textbf{False.} Take $\varphi=p$. Then $p$ is not a tautology ($w(p)=\F$
falsifies it), and $\lnot p$ is not a tautology either ($w(p)=\T$ gives $\lnot p=\F$).
So ``$\varphi$ not a tautology'' does not make $\lnot\varphi$ one: $p$ is merely
\emph{not forced true}, which is weaker than $\lnot p$ being \emph{forced true}.

\solpart{e} \textbf{True.} By Exercise 2.26, $\lnot\varphi$ is a contradiction iff
$\lnot\lnot\varphi$ is a tautology iff $\varphi$ is a tautology. Contrapositively,
if $\varphi$ is not a tautology then $\lnot\varphi$ is not a contradiction.

\solpart{f} \textbf{False.} Take $\varphi=p$, $\psi=\lnot p$: neither is a tautology
($p$ fails at $w(p)=\F$; $\lnot p$ fails at $w(p)=\T$), yet $(\varphi\lor\psi)=
(p\lor\lnot p)$ is a tautology.

\solpart{g} \textbf{True.} For every $w$, $w((\varphi\land\psi))=w(\varphi)\land
w(\psi)$; the $\land$-table yields $\T$ only when both conjuncts are $\T$. So if
$(\varphi\land\psi)$ is a tautology, $w(\varphi)=\T$ for every $w$, i.e.\ $\varphi$
is a tautology (and so is $\psi$).

\solpart{h} \textbf{True.} ``$\varphi$ is a contradiction'' means $w(\varphi)=\F$
for every $w$; its negation is that some $w$ has $w(\varphi)\neq\F$, hence (values
being only $\T,\F$) $w(\varphi)=\T$. So a non-contradiction is made true by some
assignment.

\prob{12}{Around modus ponens: Exercise 2.27 revisited}

\solpart{a} \textbf{True.} For every $w$, $w(\varphi)=\T$ and $w(\psi)=\T$ (both
tautologies), so $w((\varphi\land\psi))=\T\land\T=\T$.

\solpart{b} \textbf{True} (this is Problem 11(g), both conjuncts). For every $w$,
$w(\varphi)\land w(\psi)=\T$ forces $w(\varphi)=w(\psi)=\T$; so both are tautologies.

\solpart{c} \textbf{False.} Take $\varphi=p$, $\psi=\lnot p$. Then $(\varphi\lor
\psi)=(p\lor\lnot p)$ is a tautology, but $\varphi=p$ is not ($w(p)=\F$) and
$\psi=\lnot p$ is not ($w(p)=\T$). The disjunction is \emph{forced} true, yet
neither disjunct is forced true --- the $2.27(c)$-style trap.

\solpart{d} \textbf{True.} For every $w$: $\psi$ a contradiction gives $w(\psi)=\F$,
and $(\varphi\to\psi)$ a tautology gives $w((\varphi\to\psi))=\T$. In the
$\to$-table a value $\T$ with consequent $\F$ forces the antecedent $\F$; so
$w(\varphi)=\F$. Thus $\varphi$ is a contradiction.

\solpart{e} \textbf{True.} For every $w$: $w(\varphi)=\T$ ($\varphi$ a tautology)
and $w(\varphi)=w(\psi)$ (from $(\varphi\leftrightarrow\psi)$ a tautology), so
$w(\psi)=\T$. Hence $\psi$ is a tautology.

\solpart{f} \textbf{False.} Take $\varphi=\psi=p$. Then $(\varphi\leftrightarrow
\psi)=(p\leftrightarrow p)$ is a tautology ($w(p)\leftrightarrow w(p)=\T$ always),
but $\varphi=p$ is not a tautology ($w(p)=\F$). The biconditional is forced true
because the two sides are \emph{identical}, not because either is a tautology.

\prob{13}{Forced, impossible, or independent}

\solpart{a} \textbf{(I): $H$ forces $C$.} Suppose $\psi$ is a tautology. For every
$w$, $w(\psi)=\T$, so $w((\varphi\to\psi))=w(\varphi)\to\T=\T$. Hence
$(\varphi\to\psi)$ is a tautology --- for \emph{every} $\varphi$. Every instance of
$H$ satisfies $C$.

\solpart{b} \textbf{(III): $C$ is independent of $H$.} (This is Exercise 2.27(c).)
We give two instances of $H$ --- both $(\varphi\to\psi)$ and $\psi$ tautologies.
\begin{itemize}[nosep]
\item \emph{$C$ holds.} $\varphi=(p\lor\lnot p)$, $\psi=(q\lor\lnot q)$. Check $H$:
$\psi$ is a tautology; and $(\varphi\to\psi)$ is a tautology since its consequent
$\psi$ is (for every $w$, $w((\varphi\to\psi))=w(\varphi)\to\T=\T$). Check $C$:
$\varphi=(p\lor\lnot p)$ is a tautology. So this instance satisfies $C$.
\item \emph{$C$ fails.} $\varphi=p$, $\psi=(q\lor\lnot q)$. Check $H$: $\psi$ is a
tautology; and $(\varphi\to\psi)=(p\to(q\lor\lnot q))$ is a tautology (consequent
always $\T$). Check $C$: $\varphi=p$ is \emph{not} a tautology, since $w(p)=\F$
gives $w(p)=\F$. So this instance fails $C$.
\end{itemize}
Both satisfy $H$; one satisfies $C$ and one does not, so we are in case (III). The
lesson: $(\varphi\to\psi)$ a tautology says only $w(\varphi)\to w(\psi)=\T$ for
every $w$; when $w(\psi)=\T$ (which holds, as $\psi$ is a tautology) this is
satisfied \emph{whatever} $w(\varphi)$ is, so $\varphi$ is left completely free.
Object-level modus ponens fails to transfer to ``all assignments'' here.

\solpart{c} \textbf{(I): $H$ forces $C$.} Suppose $\varphi$ is a contradiction. For
every $w$, $w(\varphi)=\F$, so $w((\varphi\to\psi))=\F\to w(\psi)=\T$. Hence
$(\varphi\to\psi)$ is a tautology, for every $\psi$. (\emph{Ex falso}: a false
antecedent makes the implication vacuously true.)

\solpart{d} \textbf{(II): $H$ forces the failure of $C$.} Suppose $\varphi$ is a
contradiction. For every $w$, $w(\varphi)=\F$, so $w((\varphi\land\psi))=\F\land
w(\psi)=\F$. Thus $(\varphi\land\psi)$ is false under every assignment; since at
least one assignment exists (part (g)) and no formula is both a tautology and a
contradiction (part (g)), $(\varphi\land\psi)$ is \emph{not} a tautology. So no
instance of $H$ satisfies $C$. (Same $H$ as (c), opposite outcome: $\lnot$ makes
$\to$ vacuously true but $\land$ vacuously false.)

\solpart{e} \textbf{(II): $H$ forces the failure of $C$.} If $(\varphi\to\psi)$ is a
contradiction then for every $w$, $w((\varphi\to\psi))=\F$; the only $\to$-row with
value $\F$ is antecedent $\T$, consequent $\F$, so $w(\varphi)=\T$ and $w(\psi)=\F$
for every $w$. In particular $\psi$ is a contradiction, hence (by (g), and since an
assignment exists) not a tautology. So no instance of $H$ satisfies $C$. (Indeed
$H$ here is very restrictive: it forces $\varphi$ to be a tautology and $\psi$ a
contradiction.)

\solpart{f} \textbf{(III): $C$ is independent of $H$.} (This upgrades Problem
11(d).) Two instances of $H$ --- $\varphi$ not a tautology.
\begin{itemize}[nosep]
\item \emph{$C$ holds.} $\varphi=(p\land\lnot p)$, a contradiction, hence not a
tautology (it is false under, e.g., $w(p)=\T$). Then $\lnot\varphi=\lnot(p\land
\lnot p)$ is a tautology by Exercise 2.26. So $C$ holds.
\item \emph{$C$ fails.} $\varphi=p$: not a tautology ($w(p)=\F$). And $\lnot\varphi
=\lnot p$ is not a tautology ($w(p)=\T$ gives $\lnot p=\F$). So $C$ fails.
\end{itemize}
Both satisfy $H$; $C$ goes both ways, so (III). (``$\varphi$ not forced true'' is
compatible both with $\varphi$ being contingent --- then $\lnot\varphi$ is also
contingent --- and with $\varphi$ being a contradiction --- then $\lnot\varphi$ is
a tautology. This is exactly why (III), not (II).)

\solpart{g} \textbf{No formula is both a tautology and a contradiction.} Suppose
$\varphi$ were both. As a tautology, $w(\varphi)=\T$ for every truth assignment $w$;
as a contradiction, $w(\varphi)=\F$ for every $w$. It remains to know that
\emph{there is} at least one truth assignment: take any function $v\colon P\to\TF$
(for instance the constant $\T$), and by the \textbf{existence half of Theorem 2.2}
it extends to a truth assignment $\ext{v}$. For that $\ext{v}$ we get
$\T=\ext{v}(\varphi)=\F$, a contradiction. Hence no $\varphi$ is both. The proof
uses the \emph{existence} half: were the space of assignments empty, ``true under
all'' and ``false under all'' could hold vacuously at once, and the two notions
would not be incompatible.

\prob{14}{Characterizations: Exercise 2.31 revisited}

\solpart{a} \textbf{Given $\varphi$ a tautology: $(\varphi\to\psi)$ is a tautology
iff $\psi$ is a tautology.}
($\Leftarrow$) If $\psi$ is a tautology then $w((\varphi\to\psi))=w(\varphi)\to\T=
\T$ for every $w$ (this direction needs nothing about $\varphi$).
($\Rightarrow$) Suppose $(\varphi\to\psi)$ and $\varphi$ are tautologies. For every
$w$, $w(\varphi)=\T$ and $w((\varphi\to\psi))=\T$; the $\to$-row with antecedent
$\T$ and value $\T$ has consequent $\T$, so $w(\psi)=\T$. Thus $\psi$ is a tautology.

\solpart{b} \textbf{Given $\psi$ a tautology: $(\varphi\to\psi)$ is \emph{never} a
contradiction.} If $\psi$ is a tautology then $w((\varphi\to\psi))=w(\varphi)\to\T=
\T$ for every $w$, so $(\varphi\to\psi)$ is itself a tautology; being a tautology
(and, by 13(g), not simultaneously a contradiction) it is not a contradiction, for
any $\varphi$. Circumstances: none.

\solpart{c} \textbf{Refuted: $\varphi,\psi$ need not be contradictions.} Take
$\varphi=p$, $\psi=\lnot p$. Then $(\varphi\land\psi)=(p\land\lnot p)$ is a
contradiction ($w(p)\land\lnot w(p)=\F$ always), but $\varphi=p$ is not ($w(p)=\T$
gives $p=\T$) and $\psi=\lnot p$ is not ($w(p)=\F$ gives $\lnot p=\T$). (Dual to
12(c): the conjunction forced false does not force either conjunct false.)

\solpart{d} \textbf{$(\varphi\lor\psi)$ is a contradiction iff both $\varphi$ and
$\psi$ are contradictions.}
($\Rightarrow$) If $(\varphi\lor\psi)$ is a contradiction then $w((\varphi\lor\psi))
=\F$ for every $w$; the only $\lor$-row with value $\F$ has both disjuncts $\F$, so
$w(\varphi)=\F$ and $w(\psi)=\F$ for every $w$ --- both are contradictions.
($\Leftarrow$) If both are contradictions then $w(\varphi)=w(\psi)=\F$, so
$w((\varphi\lor\psi))=\F\lor\F=\F$, for every $w$. (Contrast (c): for $\lor$ the
condition really does descend to each disjunct.)

\solpart{e} \textbf{Refuted}, then characterized. Take $\varphi=p$, $\psi=\lnot p$:
$(\varphi\leftrightarrow\psi)=(p\leftrightarrow\lnot p)$ is a contradiction ($w(p)
\leftrightarrow\lnot w(p)=\F$ always), but neither $p$ nor $\lnot p$ is a
contradiction. \emph{Characterization:} $(\varphi\leftrightarrow\psi)$ is a
contradiction iff $w(\varphi)\neq w(\psi)$ for every truth assignment $w$ (i.e.\
$\varphi$ and $\psi$ take opposite values everywhere; equivalently $\psi\equiv
\lnot\varphi$). Indeed the $\leftrightarrow$-table gives $w((\varphi\leftrightarrow
\psi))=\F$ exactly when $w(\varphi)\neq w(\psi)$; so $(\varphi\leftrightarrow\psi)$
is false under \emph{every} $w$ iff $w(\varphi)\neq w(\psi)$ for every $w$.

\prob{15}{Majority rule: Exercise 2.29 revisited}

For a formula $\varphi$ write $T(\varphi)=\{v\in S:v(\varphi)=\T\}$; then
$|T(\varphi)|+|T(\lnot\varphi)|=|S|$, and $\varphi\in D$ iff $|T(\varphi)|>|S|/2$.

\solpart{a} For each $v\in S$, $v(\lnot\varphi)=\lnot v(\varphi)$, so $v\in
T(\varphi)$ iff $v\notin T(\lnot\varphi)$; hence $|T(\lnot\varphi)|=|S|-|T(\varphi)|$.
Because $|S|$ is odd, $|T(\varphi)|\neq|S|/2$ (no tie), so exactly one of
$|T(\varphi)|>|S|/2$ and $|T(\lnot\varphi)|=|S|-|T(\varphi)|>|S|/2$ holds. Thus
exactly one of $\varphi,\lnot\varphi$ is in $D$.

\solpart{b} If $(\varphi\to\theta)$ is a tautology then every $v\in S$ has
$v((\varphi\to\theta))=\T$, i.e.\ $v(\varphi)=\T\Rightarrow v(\theta)=\T$; hence
$T(\varphi)\subseteq T(\theta)$. So $|T(\theta)|\ge|T(\varphi)|>|S|/2$, giving
$\theta\in D$. (Tautological implication can only enlarge the true-set.)

\solpart{c} Take the three assignments on $\{p,q\}$
\[
v_1:(p,q)=(\T,\T),\qquad v_2:(p,q)=(\T,\F),\qquad v_3:(p,q)=(\F,\F).
\]
Majority means at least $2$ of $3$.
\begin{itemize}[nosep]
\item $p\in D$: $p$ is true under $v_1,v_2$ --- $2$ of $3$. \checkmark
\item $(p\to q)\in D$: $(p\to q)$ is true under $v_1$ ($\T\to\T$) and $v_3$
($\F\to\F=\T$), false under $v_2$ ($\T\to\F$) --- $2$ of $3$. \checkmark
\item $q\notin D$: $q$ is true only under $v_1$ --- $1$ of $3$, not a majority.
\checkmark
\end{itemize}
So $p,(p\to q)\in D$ but $q\notin D$: majority rule is not closed under modus ponens.

\solpart{d} \textbf{Refuted.} Take $\varphi=p$, $\psi=q$ and the three assignments
\[
v_1:(\T,\F),\qquad v_2:(\F,\T),\qquad v_3:(\T,\T).
\]
Then $p$ is true under $v_1,v_3$ ($2$ of $3$), so $p\in D$; $q$ is true under
$v_2,v_3$ ($2$ of $3$), so $q\in D$; but $(p\land q)$ is true only under $v_3$
($1$ of $3$), so $(p\land q)\notin D$. Hence $D$ need not be closed under $\land$.

\solpart{e} Drop oddness. Take $|S|=2$ on $\{p\}$: $v_1(p)=\T$, $v_2(p)=\F$. For
$\varphi=p$, $|T(p)|=1$, which is not $>|S|/2=1$, so $p\notin D$; likewise
$|T(\lnot p)|=1$, so $\lnot p\notin D$. Neither $p$ nor $\lnot p$ is in $D$, so the
conclusion of (a) fails (an even split is now possible).

\solpart{f} \textbf{$|S|=1$: $D$ is closed under modus ponens.} Write $S=\{v\}$. A
``majority'' of a one-element set is that one element, so $D=\{\varphi:v(\varphi)=
\T\}$. If $\varphi\in D$ and $(\varphi\to\theta)\in D$ then $v(\varphi)=\T$ and
$v((\varphi\to\theta))=\T$, i.e.\ $v(\varphi)\to v(\theta)=\T$; with $v(\varphi)=\T$
the $\to$-table forces $v(\theta)=\T$, so $\theta\in D$.

\emph{The step that breaks at $|S|=3$.} For a single assignment, ``$\varphi\in D$''
and ``$(\varphi\to\theta)\in D$'' are witnessed by the \emph{same} $v$, which then
also witnesses $\theta$: one assignment does all the work. For $|S|=3$, membership
$\varphi\in D$ is carried by a majority $M_\varphi$ (size $\ge 2$) and
$(\varphi\to\theta)\in D$ by a possibly \emph{different} majority
$M_{\varphi\to\theta}$ (size $\ge 2$). Modus ponens applies only to assignments in
$M_\varphi\cap M_{\varphi\to\theta}$ --- those that make both $\varphi$ and
$\varphi\to\theta$ true, and hence $\theta$ true --- but two $2$-element majorities
of a $3$-element set can meet in just \emph{one} assignment, which is not a
majority. That is exactly what (c) exploits: $M_p=\{v_1,v_2\}$ and $M_{p\to q}=
\{v_1,v_3\}$ meet only in $\{v_1\}$, so only $v_1$ is forced to make $q$ true ---
$1$ of $3$, not enough for $q\in D$.

\section*{Part V --- Capstone}

\prob{16}{$\star$ Repetition-free formulas: Exercise 2.30 generalized}

\solpart{a} \textbf{Gluing Lemma.} Let $\theta,\psi$ share no variable, and let
$u,w$ be truth assignments (say on a common $\Form(P,S)$ with $\vars(\theta)\cup
\vars(\psi)\subseteq P$). Define $x$ on the variables by
\[
x(r)=\begin{cases}
u(r), & r\in\vars(\theta),\\
w(r), & r\in\vars(\psi),\\
\T, & \text{otherwise.}
\end{cases}
\]
This is well-defined because $\vars(\theta)\cap\vars(\psi)=\varnothing$, so the
first two clauses never conflict. By the \emph{existence half of Theorem 2.2}, $x$
extends to a truth assignment $\ext{x}$ on $\Form(P,S)$. Now $\ext{x}$ agrees with
$u$ on every variable of $\theta$, so by \emph{Problem 8(a)} (relevance),
$\ext{x}(\theta)=u(\theta)$; likewise $\ext{x}$ agrees with $w$ on every variable
of $\psi$, so $\ext{x}(\psi)=w(\psi)$. Renaming $\ext{x}$ as $x$, this is the
required assignment. \hfill$\square$

\solpart{b} \textbf{Claim.} Every repetition-free formula of $\Form(P,S)$ is made
true by some assignment and false by some assignment. Induct on the length; IH:
this holds for all repetition-free formulas of length $\le n$.

\emph{Base ($\chi=p$).} A variable is repetition-free; $p\mapsto\T$ makes it true
and $p\mapsto\F$ makes it false (both extend to assignments by Theorem 2.2).

\emph{Step ($\chi$ repetition-free of length $n+1$).} First record two facts.
(i) The immediate subformulas of $\chi$ are repetition-free: their variables form a
subset of $\vars(\chi)$, and no variable repeats in $\chi$, so none repeats in a
subformula. (ii) In the binary case $\chi=(\theta\bin\psi)$, $\vars(\theta)\cap
\vars(\psi)=\varnothing$: a common variable would occur once in $\theta$ and once
in $\psi$, hence at least twice in $\chi$, contradicting repetition-freeness. Also
$\theta,\psi$ have length $\le n$.
\begin{itemize}[nosep]
\item $\chi=\lnot\theta$. By the IH there are assignments $a,b$ with $a(\theta)=\T$
and $b(\theta)=\F$; then $a(\lnot\theta)=\F$ and $b(\lnot\theta)=\T$, so $\lnot
\theta$ takes both values.
\item $\chi=(\theta\bin\psi)$, $\bin\in\{\land,\lor,\to,\leftrightarrow\}$. By (ii)
$\vars(\theta)\cap\vars(\psi)=\varnothing$, so the Gluing Lemma (a) applies: for
\emph{any} target pair $(s,t)\in\TF\times\TF$ there is an assignment $x$ with
$x(\theta)=s$ and $x(\psi)=t$ (take $u$ realizing $\theta=s$ and $w$ realizing
$\psi=t$, available by the IH, and glue). Now the shared feature of all five
connectives' truth tables is that \textbf{each output column is non-constant --- it
contains at least one $\T$ and at least one $\F$}:
\[
\begin{array}{c|cccc}
 & (\T,\T) & (\T,\F) & (\F,\T) & (\F,\F)\\\hline
\land & \T & \F & \F & \F\\
\lor & \T & \T & \T & \F\\
\to & \T & \F & \T & \T\\
\leftrightarrow & \T & \F & \F & \T
\end{array}
\]
So some pair $(s,t)$ makes $(\theta\bin\psi)$ true and some other pair makes it
false; gluing realizes each, giving assignments that make $\chi$ true and false.
\end{itemize}
This closes the induction. \hfill$\square$

\solpart{c} \emph{Exercise 2.30.} A formula built from $\lnot$ and $\to$ with no
repeated variable is repetition-free and uses connectives in $S$, so by (b) it is
false under some assignment; hence it is not a tautology. \emph{Stronger conclusion
delivered by the proof:} (b) shows such a formula is also true under some
assignment, so it is neither a tautology nor a contradiction --- it is
\emph{contingent}; and this holds for repetition-free formulas over the \emph{full}
set $S=\Sfive$, not only the $\{\lnot,\to\}$ fragment. (A length-$0$ repetition-free
formula, a single variable, is likewise contingent.)

\solpart{d} Take $\chi=(p\to p)$, which repeats $p$. It is a tautology, so the
conclusion of (b) --- false under some assignment --- fails for it. The step that
breaks is the binary case: here $\theta=\psi=p$, so $\vars(\theta)\cap\vars(\psi)=
\{p\}\neq\varnothing$, and the disjointness on which (b) relied does not hold.
Consequently the Gluing Lemma does not apply, and the pair $(s,t)=(\T,\F)$ needed
to make $(p\to p)$ false is \emph{unrealizable}: it would require an assignment with
$x(p)=\T$ and $x(p)=\F$ at once. Only $(\T,\T)$ and $(\F,\F)$ occur, and $\to$ is
true on both --- hence the tautology. The precise failing move is ``$\vars(\theta)
\cap\vars(\psi)=\varnothing$, so all four $(s,t)$ are realizable.''

\solpart{e} The freedom used in (a) --- setting $x(\theta)$ and $x(\psi)$ to any
pair of values at once --- is available in (b) precisely because
repetition-freeness makes $\vars(\theta)$ and $\vars(\psi)$ disjoint, so $\theta$
and $\psi$ act like independent variables. In Problem 13(b) (Exercise 2.27(c))
there is no such guarantee: $\varphi,\psi$ are arbitrary metavariables, and the map
$w\mapsto(w(\varphi),w(\psi))$ may reach only some of the four value-pairs. Distinct
metavariables are \emph{not} distinct propositional variables: $\varphi$ and $\psi$
can share variables, or be logically related (e.g.\ $\psi=\lnot\varphi$, as in
6(d)), or one can be pinned (e.g.\ $\psi$ a tautology, as in 13(b)) --- their values
are computed downstream and may be correlated. When $\psi$ is forced true, the pair
always has second coordinate $\T$, so the combination that would make $\varphi$ a
tautology is simply never demanded, and $\varphi$ stays free. The analogy with
independent variables is restored exactly when $\vars(\varphi)\cap\vars(\psi)=
\varnothing$ --- or, more generally, when $w\mapsto(w(\varphi),w(\psi))$ is onto
$\TF\times\TF$; that surjectivity is what gluing supplies for repetition-free
formulas and what is missing in general.

\bigskip
\hrule
\smallskip
\noindent\emph{Three threads, one theorem.} The uniqueness half of Theorem 2.2
drives every ``agreement'' result (1(e,g), 8, 9(b), and the determinacy in 5); the
existence half is what lets you \emph{prescribe} values --- freely on the variables
(2(d), 6, 7), across disjoint variable sets (16(a,b)), and even just to know an
assignment exists at all (13(g)). The counterexample work (11, 12(c,f), 13(b,f),
14(c,e), 15(c,d)) is the discipline of not importing the independence of distinct
variables into arbitrary metavariables --- the single idea that 16(e) makes precise.

\end{document}